% Working draft for Empirical Software Engineering (Springer).
% Before submission, move to the journal's LaTeX template (sn-jnl.cls,
% "Springer Nature LaTeX template"); the section structure already
% follows EMSE's usual research-article shape.
% Every number in Sections 4-5 comes from analysis.py on the frozen run;
% \result{} marks a value that does not exist yet. Do not fill one by hand.
\documentclass[11pt,a4paper]{article}

\usepackage[utf8]{inputenc}
\usepackage[T1]{fontenc}
\usepackage{lmodern}
\usepackage[margin=1in]{geometry}
\usepackage{microtype}
\usepackage{booktabs}
\usepackage{tabularx}
\usepackage{enumitem}
\usepackage{xcolor}
\usepackage[hidelinks]{hyperref}
\usepackage{url}

\newcommand{\result}[1]{\textcolor{red}{[RESULT: #1]}}
\newcommand{\code}[1]{\texttt{#1}}

\title{Moved, Archived, Gone: How Often Widely Used Packages Lose Their
Upstream, and Whether Dependency Tools Can Tell}

\author{Mehmet Kış\\
Ikarus AI Labs, Cambridge, United Kingdom\\
\href{mailto:mehmet@ikaruslabs.net}{\texttt{mehmet@ikaruslabs.net}}}

\date{Draft, \today}

\begin{document}
\maketitle

\begin{abstract}
\noindent\textbf{Context.} Dependency tools decide whether a dependency
is safe mostly from vulnerability advisories and lifecycle catalogues.
Changes to the dependency's upstream, such as an archived, relocated or
deleted source repository, need not appear in either.
\textbf{Objective.} We measure how often the most depended-upon packages
of npm, PyPI and Go have such a change, how many dependents they have,
and whether the change is visible through a machine-readable signal a
dependency tool could act on.
\textbf{Method.} A pre-specified cross-sectional study of
\result{N} packages, the top 1,000 and a random 1,000 from ranks
1,001--20,000 per ecosystem, measuring each package's published source
repository through the GitHub API and its registry, advisory and
lifecycle records through their primary sources.
\textbf{Results.} \result{RQ1--RQ4 headline figures}.
\textbf{Conclusion.} \result{one sentence, written after the results}.
\end{abstract}

\section{Introduction}
\label{sec:intro}

% Motivation: advisories (OSV, GHSA) and lifecycle data (endoflife.date)
% cover vulnerabilities and announced support windows. Repository-level
% changes - archival, transfer, deletion - change who can fix the code and
% where it lives, without necessarily producing either kind of record.
% Examples from TR-2026-01 (etcd relocation, Bitnami withdrawal) are
% anecdotes; this paper measures the phenomenon.

\paragraph{Research questions.}
\begin{enumerate}[leftmargin=*,label=RQ\arabic*.]
  \item \emph{Prevalence.} Among the most depended-upon packages of npm,
    PyPI and Go, what share have a source repository that is archived,
    relocated, or no longer reachable?
  \item \emph{Exposure.} How many dependent packages and repositories do
    those packages have?
  \item \emph{Signal coverage.} For packages whose source repository is
    archived or unreachable, what share carry a registry deprecation
    flag, an OSV advisory, or an endoflife.date entry?
  \item \emph{Stale references.} For relocated repositories, what share of
    registry metadata still names the old location?
\end{enumerate}

\paragraph{Contributions.}
\begin{enumerate}[leftmargin=*]
  \item A measurement of upstream source changes across three
    ecosystems, separated into strata so that it is not read as
    ecosystem-wide.
  \item A measurement of how much of it is visible to tools that read
    registries, advisories and lifecycle data.
  \item A replication package: the frozen sample, every raw response,
    the collector and the analysis.
\end{enumerate}

\section{Background and Related Work}
\label{sec:related}
% Abandonment and deprecation: Valiev et al. 2018; Avelino et al. 2019;
% Cogo et al. 2022 (npm deprecation); Miller et al. 2023.
% Advisory and lifecycle data: OSV schema, GHSA, endoflife.date.
% Supply-chain metadata: SBOM accuracy, package identity (purl),
% repository provenance (Sigstore, SLSA source provenance).
% OpenPulse / TR-2026-01: the motivating instrument; this study does not
% evaluate it.

\section{Study Design}
\label{sec:design}
% Summarise papers/prevalence/PROTOCOL.md: frame, strata and seed,
% exclusions, measures table, validation, analysis. State the freeze tag
% and the SHA-256 of the protocol, and the disclosure in its section 8
% (summary counts viewed before writing it).

\section{Results}
\label{sec:results}
\subsection{RQ1: Prevalence}
\result{per ecosystem and stratum: archived, relocated, unreachable, with Wilson 95\% intervals; excluded counts}
\subsection{RQ2: Exposure}
\result{dependent counts; median and total, with the overlap caveat}
\subsection{RQ3: Signal coverage}
\result{2x3 table: state by registry / advisory / lifecycle, and any signal}
\subsection{RQ4: Stale references}
\result{share of relocated repositories whose registry metadata names the old location}
\subsection{Validation}
\result{frame agreement; rater agreement (kappa) and which raters were independent; archived-repository classification}

\section{Discussion}
\label{sec:discussion}
% Implications for dependency tools and SBOM producers; for registries
% (publishing repository state); for agents choosing dependencies.

\section{Threats to Validity}
\label{sec:threats}
\paragraph{Construct validity.} Archived is not abandoned; unreachable
includes private repositories; lexical signal matching (estimated in the
validation sample).
\paragraph{Internal validity.} The study is descriptive and draws no
causal conclusion. Ecosyste.ms decides only the frame and the ranking;
the frame-versus-source agreement is reported.
\paragraph{External validity.} Three ecosystems, GitHub-hosted sources
only, a ranked frame; the \emph{Tail} stratum limits but does not remove
the bias towards popular packages.
\paragraph{Conclusion validity.} Proportions with intervals and no
hypothesis tests; overlapping dependent counts are never summed into a
count of affected projects.

\section{Conclusion}
\label{sec:conclusion}
\result{written after the results}

\section*{Data Availability}
The protocol, the rating guide, the collector, the frozen sample, every
raw response and the analysis code are in
\url{https://github.com/IkarusAILabs/OpenPulse/tree/main/papers/prevalence},
archived on Zenodo at \result{DOI}.

\section*{AI Assistance}
% State the tools and versions, what they were used for (code, drafting),
% and that the author made every design decision. Ratings by a language
% model, if any were used, are labelled and are not independent ratings.

\section*{Conflict of Interest}
% The author develops OpenPulse, which acts on upstream changes of the
% kind measured here. The study uses no OpenPulse code or data.

\end{document}
